\section{Counting the Length Bound}
\label{app:counting}

The same length formula applies to arbitrary and compiled readers.
The difference lies in how the number of word actions and the stopped
depth are bounded.

\subsection{Actions and obstructions of a compiled reader}

A live state has an upper and a lower subset of the $2v$ signed
variables. There are $2^{4v}$ such pairs in the table and one dead
state. Even pairs never reached from the starts can be used as inputs
of an action. Join live states by componentwise union and let
$\dead$ absorb joins. Both relational images commute with unions.
A joined image passes the guard exactly when both images pass it,
because a union is disjoint from the facts exactly when both
parts are. This gives preservation of joins for each letter,
and then for each word. Both $\dead$ and the pair of empty sets
are fixed under all word actions.

Each pair is a join of the singleton states of its two components.
There are $4v$ singleton states, and the action is determined by
its images on them. The number of actions is at most
\[
  (2^{4v}+1)^{4v}\ \le\ 2^{(4v+1)4v}.
\]
Also, a joined image is dead exactly when some component image is
dead. If all images are live, admission with a fixed target is
tested by an intersection with the union of lower and flipped
upper sets. So this admission is witnessed by one singleton's
image. More generally, death or admission with the fixed target
holds after a word exactly when it holds for one of the singleton
states carried by the starting pair. A word read from any starting
pair carrying that singleton then has the same obstruction, up to
death in the admission case. These are the label hypotheses of
\cref{app:charging}. Taking $J=4v$ there proves the compiled part
of \cref{lem:depth}.

\subsection{Combining the bounds}

Given a hole, \cref{lem:stopping} supplies an extension with stopped
depth bounded by $C$. By \cref{lem:periodic} and \eqref{eq:offset}, its
moving coordinates give rays of period at most $D$ and offsets.
The resulting periodic tuple can be shortened as in
\cref{lem:buffered}. Its onsets and stopped lengths are bounded by
$S$ and its periods by $D$. Then \cref{lem:truncation} gives a hole
of length at most $\LBound(k,N,K,C)$.

To estimate the formula, suppose $k,C\le a$ and $N,K\le 2^a$.
Write $t=a+16$. In the following estimates we use $k\le 2^a$,
$C+2\le t$, and $t\ge16$. The quantities $k+2$, $N+2$,
$C+N+3$, $K+1$, and $2k+3$ are all at most $2^t$, since each
is at most $2^{a+1}+3\le 2^{a+3}$. Thus
\begin{align*}
 D &\le 2^{(2t+1)t} \le 2^{3t^2},\\
 B+2 &\le 2^t D + 6 \le 2^{3t^2+t+1}\le 2^{4t^2},\\
 M &\le 1+2^{8t^2+t} \le 2^{9t^2}.
\end{align*}
For the finite bound used for $S$, its state and action arguments
are both bounded by $2^{10t^2}$. Its budget
$R=(MN+M(K+1)+1)(2k+3)^{3(C+1)}$ is at most
$2^{13t^2+2}\le 2^{14t^2}$. Hence
$Q=2R+k+2\le 2^{14t^2+2}\le 2^{15t^2}$.
Substituting in the finite bound gives
\[
       S\le 2^{15t^4+t},
       \qquad H=S+ND+1\le 2^{15t^4+t+2}.
\]
Each summand inside the parentheses of \eqref{eq:length-bound}
is at most $2^{15t^4+2t+2}$, including $kH$ and $2D$.
There are five summands. Multiplying by $k\le 2^t$ gives
\[
      \LBound(k,N,K,C) \le 2^{15t^4+3t+5}\le 2^{t^5}.
\]
The last inequality follows from $3t+5\le t^4$ and $t\ge16$.

For arbitrary readers, choose $K=N^N$ by counting all state maps
and $C=(N+1)k$ by \cref{lem:depth}. These choices meet the preceding
hypotheses with $a=(k+N+2)^2$: the depth and $k$ are at most $a$,
and $N^N\le 2^{N^2}$ since $N\le 2^N$. For compiled readers use
$N=2^{4v}+1$, $K=2^{(4v+1)4v}$, and $C=(4v+1)k$.
Take $a=(4v+k+2)^2$. The depth and $k$ are again at most $a$,
and the exponents $4v+1$ and $(4v+1)4v$ are bounded by $a$.
In both cases $(s+2)^2+16\le(s+16)^2$, with $s=k+N$ or
$s=4v+k$. Raising to the fifth power proves the two estimates
in \cref{thm:short}. All these counts use states, actions, and
depth, independently of the alphabet size.
